\documentclass{beamer}
\input{tbook-deck}
\title{Golden Integers and Observability}
\subtitle{An overview of the example monograph}
\author{Example speaker}
\institute{}
\date{}
\begin{document}
% Manuscript: Preface, scope of the example monograph
\begin{frame}
  \titlepage
  \note{Introduce the sample monograph. The examples show how a talk can cite the book's frozen results while keeping its own exposition. Replace the speaker and institution fields for a real seminar.}
\end{frame}
% Manuscript: Contents, Chapters 1 and 2
\begin{frame}{Outline}
  \tableofcontents[hideallsubsections]
  \note{The main talk distinguishes Lean verification, a written proof and an open question. Optional technical material and the verification ledger follow the appendix marker.}
\end{frame}
\section{Golden integers}
\subsection{The carry deficit}
% Manuscript: Chapter 1, Theorem 1.3
\begin{frame}{The carry deficit takes three values}
  \begin{bookresult}{thm-deficit-three-valued}
    For all natural numbers $v_1,v_2$, the normalized carry deficit satisfies
    \[
      c(v_1,v_2)\in\{-1,0,1\}.
    \]
  \end{bookresult}
  The bound holds for arbitrarily large inputs.
  \note{Recall that the deficit measures the failure of the golden-integer embedding to preserve addition. The green tag links to the pinned Lean source.}
\end{frame}
% Manuscript: Chapter 1, Proposition 1.12 and proof pf-deficit-frequencies
\begin{frame}{Limiting frequencies: a written proof}
  \begin{bookresult}{rem-areas-not-frequencies}
    For $0\le v_1,v_2<n$, let $a_n^\pm$ be the proportions with deficit $\pm1$. Then
    \[
      \lim_{n\to\infty}a_n^+=\frac{1}{2\varphi^2},\qquad
      \lim_{n\to\infty}a_n^-=\frac{1}{2\varphi^4}.
    \]
  \end{bookresult}
  Weyl equidistribution, rectangle products and grid approximations give the limits.
  \note{This is the book's written argument, derived from Weyl's theorem. It has no matching frozen Lean verification. The source DOI or repository and formalization request appear automatically; the purple badge does not count as Lean verification.}
\end{frame}
\section{Refutation}
\subsection{A floor formula}
% Manuscript: Chapter 2, Theorem 2.1
\begin{frame}{A counterexample to the floor formula}
  For a positive integer $n$, let $a(n)$ be the greatest natural number $k$ with $(1+1/n)^k\leq 2$.
  \begin{bookresult}{thm-greathouse-refuted}
    The formula
    \[
      a(n)=\left\lfloor\left(n+\tfrac12\right)\log 2\right\rfloor
    \]
    does not hold for every $n\geq1$.
  \end{bookresult}
  \note{This is the refutation result. A counterexample establishes falsity of the universal formula; it does not establish the least counterexample.}
\end{frame}
% Manuscript: Chapter 2, Open problem 2.2
\begin{frame}{The least counterexample}
  \begin{bookresult}{prob-greathouse-minimal}
    Is $n_0=1121626023352383$ the least input at which the formula fails? What characterizes the inputs where it fails?
  \end{bookresult}
  \note{Keep the distinction between the checked refutation and the open minimality question. This block has no GID and cannot receive a verified badge.}
\end{frame}
\section{Pictures}
\subsection{A carry and its geometry}
% Manuscript: Chapter 1, visualization golden-carry, card plus
\begin{frame}{A positive carry deficit}
  \centering
  \BookViz{golden-carry}{plus}
  \par\smallskip\small Click the snapshot to open the live experiment.
  \note{This paper snapshot is a still of the book visualization, with the same card and language. The illustration explains the example; the Lean result supplies the general guarantee.}
\end{frame}
% Manuscript: Chapter 1, figure phase-square
\begin{frame}{The phase square}
  \centering
  \BookFigure{phase-square}
  \par\smallskip\small The figure uses the book's print palette and labels.
  \note{The canonical SVG is regenerated by the book's figure script. The slide build materializes its paper palette and converts it to a PDF asset.}
\end{frame}
\appendix
\section{Technical appendix}
\subsection{Verification scope}
% Manuscript: Preface, verification model
\begin{frame}{The scope of verification}
  Each cited result inherits its evidence from the book.
  \begin{itemize}
    \item The speaker writes the block's mathematical explanation.
    \item The book supplies its kind, number and title.
    \item The lock supplies Lean sources; written proofs supply a DOI or source repository and request.
  \end{itemize}
  Pinned source: \PinnedCommit
  \note{The identity chain is checked mechanically; fidelity of the speaker's prose to the cited statement remains an author and reviewer responsibility.}
\end{frame}
% Manuscript: Formal verification ledger, cited results only
\FormalVerification
\end{document}
